\section{Related work}\label{sec:related}
\paragraph{Conformal prediction and its evaluation.} Split (inductive) conformal prediction calibrates a score threshold on held-out data; its coverage conditional on the calibration set follows a Beta law, and \citet{angelopoulos2021gentle} recommend comparing the empirical distribution of coverage over calibration draws with this law to separate implementation errors from fluctuation, and stress that adaptivity must be assessed separately from coverage. For classifiers, \citet{angelopoulos2021raps} introduce regularised adaptive sets and size-stratified coverage violation (SSCV), and \citet{ding2023clustered} evaluate class-conditional coverage gaps (CovGap) and set size over repeated splits when classes have few calibration points. \citet{jung2023batch} give group-conditional guarantees for any score and group collection; we do not run them (Section~\ref{sec:disc}). These works treat the softmax output as a score input and do not use the geometry of the simplex.

\paragraph{Multivariate and constrained responses.} \citet{feldman2023calibrated} calibrate multi-output quantile regions and compare region volume; \citet{principato2024hierarchical} show for responses satisfying linear equality constraints that projecting the point prediction onto the constraint set cannot increase an $A$-norm score, so regions shrink without losing coverage. Sum-to-one is such a constraint, but the simplex has additional inequality constraints $y\ge0$ that their equality-constraint result does not address.

\paragraph{Composition-valued responses.} The closest setting is \citet{javanmardi2024credal}, who conformalise credal sets: the response is an annotator label distribution and the output is a split-conformal region of the simplex defined through total variation, Wasserstein, KL, an inner-product score or a Dirichlet likelihood. Efficiency is measured as the fraction of a grid of the simplex covered (resolution $0.005$, $K=3$ in their main experiment), which does not scale to larger $K$. They report average coverage, but do not compare calibration schemes, probe conditional coverage, or use log-ratio geometry. We reuse the setting and the scores as a starting point and add the evaluation questions above.

\paragraph{Compositional data analysis.} \citet{aitchison1982statistical} models compositions through log-ratios; the Aitchison distance is the Euclidean distance between centred log-ratio (CLR) coordinates and requires strictly positive parts. Zeros, which occur whenever labels are one-hot or components fall below a detection limit, need a pre-declared replacement rule; we make that rule an experimental factor.
